\subsection{Verdict against the criteria fixed in advance}
\label{sec:verdict}

\S\ref{sec:criteria} committed to three CIFAR-10 criteria before any run, and to a
task-defined standard on the constructed tasks.

\begin{ttkey}
\lead{1. Beat the greedy stack} (\accMctm{}\%). Cleared by all three repairs.

\lead{2. Beat the random-layer-1 control} (\accMctmRandom{}\%). Cleared by one:
\ideaB's controller run over that same random layer, \randBestAcc\,$\pm$\,\randBestSd{}\%.
Over the \emph{trained} layer the best that the same controller manages is \sizeBestAcc{}\%,
short of it.

\lead{3. Beat a single layer at the same clause budget} (\accSingleMatched{}\%, or
\accSingleRf{}\% with the stack's receptive field). Cleared by the same arm, and only by it.
Every other two-layer configuration in either report sits below \accSingleRf{}\%.

\lead{And on the tasks built to answer the question,} \ideaA\ closes the headroom entirely:
\synNearAutoAuto{}\% on \code{near}, where a single layer is at chance by construction, and
\synXorAutoAuto{}\% on \code{xor}, where it is capped at $75\%$ --- both past the first
report's hand-built oracle.

\lead{\ideaC\ fails everywhere,} and for the same reason in each place: the credit rule
saturates or erases layer~1 within an epoch on CIFAR-10 and on \code{near} alike.
\end{ttkey}

\subsection{What made the difference}

Across all \nArms\ arms the thing that predicts accuracy is not which of the three ideas was
applied but what the first layer's clauses ended up \emph{looking like}: short, and firing
often enough to be conjoined. Every arm that got there did well and every arm that did not did
badly, whether it got there by an unsupervised objective, by a post-hoc controller, or by not
being trained at all.

That reading also explains the two results that looked paradoxical. The first report's
hand-built oracle lost to greedily trained clauses on \code{near} because an exact
$4\times4$ detector fires almost nowhere --- perfect features that never fire are not usable
features. And \emph{random} clauses beat trained ones on CIFAR-10 because three random literals
are short and fire often, which is the whole specification.

\subsection{What the second layer is worth, and where}

With a repaired first layer the two-layer stack is finally worth its clause budget on CIFAR-10
--- \randBestAcc{}\% against \accSingleRf{}\% for the best single layer --- and on
\code{near} it is the difference between chance and \synNearAutoAuto{}\%. The size of the
advantage tracks how much the task needs what the second layer supplies, which is a conjunction
between spatially separated parts: a few points on natural images whose classes do not
obviously require one, and everything on a task constructed so that nothing else can express
the label.

\S\ref{sec:l2learned} is the corroborating measurement, and it is the one to keep. Over a
repaired first layer a layer-2 clause on \code{near} carries \synltwoNearCalibTask\ literals
--- ``channel $a$ fires here \textsc{and} channel $b$ fires there'' --- where over the greedy
layer it carried \synltwoNearMctmTask. On CIFAR-10 the same repair takes it from \sizeMctm\ to
\sizeMctmRandomCalib. A Tsetlin machine whose second layer is a thousand-literal negation soup
has lost the property it was chosen for, and fixing the first layer's density gets it back.

\subsection{What a feature layer has to satisfy}

Across both reports, a frozen feature layer under a spatial second layer has to satisfy three
conditions at once. Each repair fixes one, and \ideaC\ breaks one.

\begin{table}[h]
\centering\small
\caption{What a layer-2 clause needs from layer~1, and what supplies it.}
\label{tab:conditions}
\begin{tabular}{@{}p{0.24\textwidth}p{0.43\textwidth}p{0.26\textwidth}@{}}
\ttoprule
\thead{condition} & \thead{why} & \thead{supplied by} \\
\ttmidrule
\lead{Dense enough to conjoin}
  & A conjunction of $j$ bits each firing with probability $p$ is satisfied with probability
    ${\sim}p^{j}$; below a fraction of a percent, layer~2 can only include negations. It is
    why the hand-built oracle lost to greedy clauses on \code{near}.
  & \ideaB\ directly; \ideaA\ as a side effect. \\[2pt]
\lead{Short, and about the right things}
  & Density alone is not enough: random clauses are dense and reach
    \synNearRandomRandom{}\% on \code{near} against \synNearAutoAuto{}\% for autoencoder
    clauses at a similar rate. On CIFAR-10 cutting each greedy clause to \sizeBestTarget\
    literals is worth more than any rate target alone.
  & \ideaB's size target; \ideaA, whose clauses come out short unasked; and a random
    initialisation, which is short by construction. \\[2pt]
\lead{Stable enough to be named}
  & Layer-2 clauses are conjunctions over \emph{channel indices}. If layer-1 clause $k$
    changes, every clause including channel $k$ becomes a rule nobody trained.
  & Freezing --- which is exactly what \ideaC\ gives up. \\
\ttbotrule
\end{tabular}
\end{table}

Greedy supervised training violates the first two, which is what both reports have been
measuring. A random layer satisfies the last two for free --- three literals long, and never
changing --- and the controller supplies the first without needing a label. That combination
is the best arm in either report, and it contains no supervised feature learning at all.

\subsection{What to do next}

\lead{Do not train the feature layer. Calibrate it.} The strongest CIFAR-10 configuration in
either report has no supervised feature learning in it: random clauses, adjusted by a
label-free controller to a target firing rate, \randBestAcc{}\% against \accSingleRf{}\% for
the best single layer. It is also the cheapest --- there is no layer-1 training to pay for.
Whether that survives deeper stacks, larger clause budgets and other datasets is the obvious
thing to check, and it is cheap to check.

\lead{If a layer \emph{is} trained, cut it down afterwards.} \S\ref{sec:sizetarget} takes a
greedily trained layer and deletes all but \sizeBestTarget\ of each clause's
\densityGreedySize\ literals, choosing what to keep by firing rate. That is worth
\accMctm{}\%\,$\to$\,\sizeBestAcc{}\% for seconds of work, and it needs no labels either.

\lead{Report the density of any feature layer, always.} One number, one forward pass, and it
explains the oracle paradox of the first report, the dose-response curves of
\S\ref{sec:density}, and the difference between a layer-2 clause with
\posMctmRandomCalib\ positive literals and one with \posMctm.

\lead{Credit propagation needs a different rule, not a tuned one.} \S\ref{sec:credit} argues
that the failure is structural --- the ratio of Type~Ia to Type~Ib events is what holds a
Tsetlin clause at a workable size, and the credit path replaces a clause's own match rate with
a layer above's. Any rule that keeps that substitution inherits the same instability. The
promising direction is to restore the missing quantity: gate credit on the layer-1 clause's
\emph{own} statistics, which \ideaB's controller already measures.

\lead{And the question both reports now sharpen.} The second layer's advantage is a few points
on CIFAR-10 and the whole task on \code{near}, and the difference between them is not scale
but whether the label depends on a relation between spatially separated parts. Object
counting, relative position, containment and same-different are where a spatial second layer
should pay, and none of them is tested here.

\subsection{Limitations}

\lead{Scale.} These are small machines by CIFAR-10 standards --- hundreds of clauses where the
literature uses thousands --- chosen so that twenty-odd arms across three seeds fit in a few
GPU-hours. The comparisons are between arms at a matched budget and that is what the
conclusions rest on; the absolute accuracies are not CIFAR-10 results for Tsetlin machines.

\lead{Two constructed tasks.} \code{near} and \code{xor} were designed to have known ceilings,
not to be representative. The claim they support is a conditional one --- when a task needs a
conjunction between separated parts, repairing layer-1 density makes the stack able to find it
--- and not that such tasks are common.

\lead{One autoencoder objective, one split.} \ideaA\ is tested with centre-from-surround
reconstruction at one specificity. A masked-bit or contrastive objective might behave
differently, and \S\ref{sec:readout} shows the global-OR readout cannot measure the
information claim directly once the layer is dense.

\lead{The credit rule is one rule.} Table~\ref{tab:creditrule} is a reading of a one-sentence
proposal, and \S\ref{sec:credit} records the three places a different reading was tried and
rejected for a stated reason. Crediting several literals per event, propagating through
negated literals, or accumulating credit across batches are untested variants.

\lead{Hyper-parameters.} $T$ is the clause count in every layer and was not tuned per arm;
layer-2 specificity is tied to layer-1's. A wider search could move the CIFAR-10 numbers; it
is unlikely to close a twelve-point gap to an untrained baseline.
